\documentclass[11pt,onecolumn]{IEEEtran}
\usepackage{amsmath,amssymb,amsthm}
\usepackage[hidelinks]{hyperref}

\newtheorem{theorem}{Theorem}
\newtheorem{proposition}[theorem]{Proposition}
\newtheorem{lemma}[theorem]{Lemma}
\newtheorem{corollary}[theorem]{Corollary}
\newtheorem{remark}[theorem]{Remark}

\newcommand{\E}{\mathbb{E}}
\newcommand{\R}{\mathbb{R}}
\newcommand{\tr}{\operatorname{tr}}
\newcommand{\PC}{P_{\mathcal C}}
\newcommand{\lmin}{\lambda_{\min}}
\newcommand{\cU}{\mathcal U}
\newcommand{\Rfix}{R_{\mathrm{fixed}}}
\newcommand{\RC}{R_{\mathcal C}}

\begin{document}

\title{The Price of a Fixed Weight When the Read Kernel Moves:\\ Fixed-Covariance
Gaussian Test Channels\\ (T1 addendum, working note)}
\author{Andrew~H.~Bond%
\thanks{Working note of the Observation Theory program, v0.5, 2026-11-22. Drafted by
Claude (Anthropic) at the author's request. The scope follows review notes of
2026-11-21 that the author passed on, which also gave the example of
Section~\ref{sec:example}. Theorems~\ref{thm:sharp} and \ref{thm:reduce}, Propositions~\ref{prop:exact}
and~\ref{prop:gapex} and the comparison with prior work follow derivations and pointers
in a review of v0.2, dated 22 November 2026, and in further review notes of the same
day, all passed on by the author; Claude re-derived the mathematics and read the cited
passages of the prior work. No statement by the author of having read the proofs is on
record. The T1 note \cite{t1} is frozen at v4.1 and is not changed by this
addendum.}\\
\emph{Senior Member, IEEE}\\
Department of Computer Engineering, San Jos\'e State University\\
\texttt{andrew.bond@sjsu.edu} \textbar\ ORCID: 0009-0003-2599-6158}
\maketitle

\begin{abstract}
The T1 note learns one quadratic weight for a consumer and serves codes built for that
weight. T1 bounds the price of that weight when the directions the consumer ignores are
the same at every source point and its sandwich assumption holds. This addendum treats a
consumer whose ignored directions change with the point. Its statements are about one
class of test channels: Gaussian backward channels whose error covariance does not depend
on the source point, which are the test channels behind the codes T1 serves. For that
class, the least rate that meets a small consumer loss is found up to a term that vanishes
with the loss. Its leading term is half the rank of the averaged weight times the
logarithm of the inverse loss. When the consumer map has a measurable right inverse, a
closed convex image and an output with a regular law, the consumer's own
rate--distortion function is that of its output, so the fixed-covariance class exceeds it
by a term that grows like the logarithm of the inverse loss, with coefficient half the
difference between the rank of the averaged weight and the output dimension. For the
example $x_1+a\sin x_2$ the loss has a closed form and an explicit part of the lower bound's
range follows. These are statements about test channels and information rates, not about every
code. What a learner needs in order to obtain the right inverse is left to a separate
item.
\end{abstract}

\section{Setting}\label{sec:setting}

Logarithms are natural and rates are in nats. The source symbol is $T\sim N(0,\Sigma)$ in
$\R^p$ with $\Sigma\succ0$. As in \cite{t1}, a consumer acts through its weighted map
$g:\R^p\to\R^k$, with the loss
\[
\ell_{\mathcal C}(x,\hat x)=\|g(\hat x)-g(x)\|^2 ,
\]
the local read metric $\PC(x)=J_g(x)^{\mathsf T}J_g(x)$, and $Q=\E\,\PC(T)$. We assume
throughout that $g$ is $C^3$ with $M_j=\sup_x\|D^jg(x)\|<\infty$ for $j=1,2,3$. Write
$q=\operatorname{rank}Q$.

\textbf{Fixed-covariance test channels.} For $0\prec\Delta\preceq\Sigma$ let
$\hat T\sim N(0,\Sigma-\Delta)$ and $Z\sim N(0,\Delta)$ be independent and $T=\hat T+Z$.
This is the test channel of the Gaussian frontier codes of the series, whatever weight
they are built for: its error covariance does not depend on the source point. Its rate
and consumer loss are
\[
r(\Delta)=\tfrac12\log\det(\Sigma\Delta^{-1}),\qquad
L(\Delta)=\E\,\ell_{\mathcal C}(T,\hat T),
\]
and we define
\[
\Rfix(D)=\inf\{\,r(\Delta):0\prec\Delta\preceq\Sigma,\ L(\Delta)\le D\,\}.
\]
This is a statement about test channels. T1's coding theorem carries a test channel's
rate and consumer loss to codes in the memoryless setting \cite[Theorem ``From test
channels to codes'']{t1}. It does not show that every code chosen from a weight is bounded
below by $\Rfix$, and this note does not claim that.

\textbf{The consumer's own rate.} Let
\[
\RC(D)=\inf\{\,I(T;\hat T):\E\,\ell_{\mathcal C}(T,\hat T)\le D\,\},
\]
the infimum over all conditional laws of $\hat T\in\R^p$ given $T$.

\textbf{What T1 supplies.} For every $0\prec\Delta\preceq\Sigma$, with $s=\tr\Delta$,
\cite[Theorem ``The test channel is exact to second order'']{t1} gives
$|L(\Delta)-\tr(Q\Delta)|\le\rho(\Delta)\le\rho_b(\Delta)$, where T1 defines and its
proof derives
\[
\rho_b(\Delta)=\bigl(\tfrac74M_2^2+2M_1M_3\bigr)s^2+\tfrac{\sqrt5}2M_2M_3s^{5/2}
+\tfrac5{12}M_3^2s^3 .
\]
Since $s\le\tr\Sigma$, this is at most $Ks^2$ with
\begin{equation}\label{eq:K}
K=\tfrac74M_2^2+2M_1M_3+\tfrac{\sqrt5}2M_2M_3(\tr\Sigma)^{1/2}+\tfrac5{12}M_3^2\tr\Sigma .
\end{equation}

\section{Kernel geometry}

The picture is this. The level sets of $g$ are the fibers the consumer cannot tell apart.
Where $J_g$ has constant rank they foliate the source space, and $\ker J_g(x)$ is the
tangent space of the fiber through $x$. The read metric $\PC(x)$ measures displacement
across the fibers at $x$ and ignores displacement along them, and $Q$ averages that
transverse metric over the source. A right inverse of $g$, as used in
Section~\ref{sec:output}, picks one point on each fiber: it is a measurable selection of
fiber representatives. To call it a transversal to the foliation would need more
regularity than measurability, and none is assumed. A fixed-covariance channel has one error shape for every point, so it cannot
put its error along the fibers when their direction turns; that is the source of the
price below.

\begin{lemma}\label{lem:kernel}
$\ker Q=\{v:J_g(T)v=0\ \text{almost surely}\}$.
\end{lemma}

\begin{proof}
$v^{\mathsf T}Qv=\E\|J_g(T)v\|^2$, which is zero exactly when $J_g(T)v=0$ almost surely.
\end{proof}

So $q$ counts the directions that the consumer reads at some source points, and
$p-q$ the directions it reads at none. If $J_g$ has rank $r$ at almost every point, then
$q\ge r$, and $q=r$ means that the kernels $\ker J_g(T)$ agree almost surely: each
contains $\ker Q$ and has the same dimension. We say the kernel \emph{moves} when
$q>r$. This is geometry only. It says nothing yet about rates.

\section{The fixed-covariance rate when $Q$ has full rank}

\begin{theorem}[Fixed-covariance channels pay for every direction]\label{thm:A}
Suppose $Q\succ0$, let $\lambda=\lmin(Q)$, let $K$ be any positive constant not less than the value in \eqref{eq:K} (the remainder
bound holds for any such constant; when the value in \eqref{eq:K} is positive we take
it), and
let $s_0=\lambda/(2K)$.
\begin{enumerate}
\item[(a)] If $0\prec\Delta\preceq\Sigma$ and $\tr\Delta\le s_0$, then
$\tr\Delta\le 2L(\Delta)/\lambda$.
\item[(b)] $D_0=\inf\{L(\Delta):0\prec\Delta\preceq\Sigma,\ \tr\Delta\ge s_0\}$ is
positive (and is $+\infty$ if $s_0>\tr\Sigma$, when the set is empty).
\item[(c)] For $D<D_0$,
\[
\Rfix(D)\ \ge\ \tfrac p2\log\frac{p\lambda}{2D}+\tfrac12\log\det\Sigma .
\]
\item[(d)] $\Rfix(D)\le\tfrac p2\log(1/D)+c$ for $D<D_1$, with
\[
c=\tfrac12\log\det(\Sigma Q)+\tfrac p2\log(2p),\qquad
D_1=2p\min\Bigl\{\lmin(\Sigma)\lambda,\ \frac{p}{K(\tr Q^{-1})^2}\Bigr\}.
\]
\end{enumerate}
Hence $\Rfix(D)=\tfrac p2\log(1/D)+O(1)$ as $D\to0$.
\end{theorem}

\begin{proof}
(a) With $s=\tr\Delta\le s_0$,
$L(\Delta)\ge\tr(Q\Delta)-Ks^2\ge\lambda s-Ks^2\ge\lambda s/2$.

(b) Let $\mathcal K=\{\Delta:0\preceq\Delta\preceq\Sigma,\ \tr\Delta\ge s_0\}$, a compact
set that contains the set in (b). If it is empty there is nothing to prove. Extend $L$ to singular $\Delta$ by the same
construction, with $\hat T=(\Sigma-\Delta)^{1/2}W_1$ and $Z=\Delta^{1/2}W_2$ for
independent standard Gaussian vectors $W_1,W_2$. The matrix square root is continuous on
positive semidefinite matrices, $g$ is continuous, and
$\ell_{\mathcal C}(\hat T+Z,\hat T)\le M_1^2\|Z\|^2\le M_1^2\|\Sigma\|\,\|W_2\|^2$, so $L$
is continuous on $\mathcal K$ by dominated convergence. It remains to show $L>0$ on
$\mathcal K$. Suppose $L(\Delta)=0$ for some $\Delta\ne0$, and let $V$ and $W$ be the
ranges of $\Delta$ and $\Sigma-\Delta$. Then $g(\hat t+z)=g(\hat t)$ for almost every
$(\hat t,z)\in W\times V$: the laws of $\hat T$ and $Z$ are nondegenerate Gaussian laws
on $W$ and $V$, so each has full support there and is equivalent to Lebesgue measure on
its subspace. A continuous function that vanishes almost everywhere on $W\times V$
vanishes everywhere on it, so the identity holds for every such pair. Because
$\Sigma\succ0$, $V+W=\R^p$. Take $x\in\R^p$ and $v\in V$, and write $x=w+v'$ with
$w\in W$, $v'\in V$. Then $g(x+v)=g(w)=g(x)$. So $g$ is constant along $V$, hence
$J_g(x)v=0$ for all $x$ and all $v\in V$, and $v^{\mathsf T}Qv=0$ for some $v\ne0$,
which contradicts $Q\succ0$. A continuous positive function on a compact set has a
positive minimum, and $D_0$ is at least that minimum.

(c) Let $D<D_0$ and $L(\Delta)\le D$. By (b), $\tr\Delta<s_0$, so by (a)
$\tr\Delta\le2D/\lambda$. By the inequality of arithmetic and geometric means,
$\det\Delta\le(\tr\Delta/p)^p\le(2D/(p\lambda))^p$, and
$r(\Delta)=\tfrac12\log\det\Sigma-\tfrac12\log\det\Delta$.

(d) Take $\Delta_t=tQ^{-1}$. For $t\le\lmin(\Sigma)\lambda$ we have
$\Delta_t\preceq\Sigma$. Then $\tr(Q\Delta_t)=pt$ and
$L(\Delta_t)\le pt+K(\tr Q^{-1})^2t^2\le 2pt$ once
$t\le p/(K(\tr Q^{-1})^2)$. With $t=D/(2p)$ the loss is at most $D$ and
$r(\Delta_t)=\tfrac12\log\det(\Sigma Q)+\tfrac p2\log(2p/D)$.
\end{proof}

\begin{remark}[Scope of the theorem]\label{rem:scope}
The theorem needs $Q\succ0$. The case $q<p$ is reduced to it by
Theorem~\ref{thm:reduce}, which has $q$ in place of $p$. The constant $D_0$ comes from a compactness
argument, and in general we have no explicit lower bound for it;
Proposition~\ref{prop:exact} gives one for the example. It does have an explicit upper
bound, for instance when $s_0<\lmin(\Sigma)$ (a sufficient condition, not a necessary
one). Let $v$ be a unit eigenvector of $Q$ for $\lambda$ and
$\Delta_\varepsilon=s_0vv^{\mathsf T}+\varepsilon I$. For small $\varepsilon>0$ this is
admissible with $\tr\Delta_\varepsilon\ge s_0$, and
$L(\Delta_\varepsilon)\le\tr(Q\Delta_\varepsilon)+K(\tr\Delta_\varepsilon)^2$, so letting
$\varepsilon\to0$,
\[
D_0\ \le\ \lambda s_0+Ks_0^2=\frac{3\lambda^2}{4K}.
\]
Then the range $D<D_0$ for which part (c) is proved lies below $3\lambda^2/(4K)$, and we
do not show how far below. This limits the proof's range; it does not say the inequality
of part (c) fails for larger $D$.
The bound is over fixed-covariance Gaussian test channels only.
\end{remark}

\section{The consumer's rate as an output rate}\label{sec:output}

Let $U=g(T)\in\R^k$ and $\cU=g(\R^p)$. The image $\cU$ is a continuous image of a
$\sigma$-compact set, so it is $\sigma$-compact and Borel. For a Borel set
$\mathcal A\subseteq\R^k$ containing $\cU$, let
\[
R_U^{\mathcal A}(D)=\inf\{\,I(U;\hat U):\hat U\in\mathcal A\ \text{almost surely},\
\E\|\hat U-U\|^2\le D\,\},
\]
the infimum over conditional laws of $\hat U$ given $U$, and $R_U=R_U^{\R^k}$. Like
$\RC$, these are informational quantities. This section invokes no coding theorem. To
carry its conclusions to codes through T1's theorem one would also need
$\E\|s(\hat U)\|^2<\infty$ for the section $s$ below, which a general measurable section
does not give.

\begin{proposition}[A section turns the consumer's rate into an output rate]\label{prop:B}
\begin{enumerate}
\item[(a)] $\RC(D)\ge R_U^{\cU}(D)\ge R_U(D)$.
\item[(b)] If there is a measurable $s:\cU\to\R^p$ with $g(s(u))=u$ for all
$u\in\cU$, then $\RC(D)=R_U^{\cU}(D)$.
\item[(c)] If moreover $\cU$ is closed and convex, then $\RC(D)=R_U(D)$.
\end{enumerate}
\end{proposition}

\begin{proof}
(a) For any admissible $\hat T$, the variable $\hat U=g(\hat T)$ takes values in $\cU$
and $\E\|\hat U-U\|^2=\E\,\ell_{\mathcal C}(T,\hat T)\le D$. Since $U$ is a function of
$T$ and $\hat U$ of $\hat T$, $I(T;\hat T)\ge I(U;\hat U)\ge R_U^{\cU}(D)$. The second
inequality holds because $\cU\subseteq\R^k$.

(b) Take any conditional law of $\hat U\in\cU$ given $U$ with
$\E\|\hat U-U\|^2\le D$, apply it to $U=g(T)$, and set $\hat T=s(\hat U)$. Then
$\ell_{\mathcal C}(T,\hat T)=\|\hat U-U\|^2$, and $T\to U\to\hat U\to\hat T$ is a Markov
chain, so $I(T;\hat T)\le I(U;\hat U)$. Taking the infimum gives
$\RC(D)\le R_U^{\cU}(D)$.

(c) Let $\pi$ be the Euclidean projection onto $\cU$. It does not increase the distance
to any point of $\cU$, and $U\in\cU$, so for any $\hat U\in\R^k$,
$\|\pi(\hat U)-U\|\le\|\hat U-U\|$, while $I(U;\pi(\hat U))\le I(U;\hat U)$. Hence
$R_U^{\cU}(D)\le R_U(D)$.
\end{proof}

When the image is curved and not convex, (b) still applies and the reproductions must
stay on the image. Part (c) is what lets unrestricted Euclidean reproductions be used.

The proposition shows that describing the output suffices for a reproduction map that
uses a section. We call the holder of that map the decoder, informally: no code is
constructed here. The substance of hypothesis (b) is that the decoder holds a section,
not that one exists. It does not show that a learner must learn the level sets of $g$. That would
need a lower bound of its own.

\begin{corollary}[The rate gap]\label{cor:C}
Suppose the hypotheses of Theorem~\ref{thm:A} and of Proposition~\ref{prop:B}(c) hold,
that $U$ has a density with finite differential entropy $h(U)$, and that its covariance
$K_U$ is finite and positive definite with least eigenvalue $\kappa$. Then for
$D<\min\{D_0,D_1,k\kappa\}$,
\[
\tfrac{p-k}2\log\frac1D+c_-\ \le\ \Rfix(D)-\RC(D)\ \le\ \tfrac{p-k}2\log\frac1D+c_+ ,
\]
with
$c_-=\tfrac p2\log\tfrac{p\lambda}2+\tfrac12\log\det\Sigma-\tfrac12\log\det K_U-\tfrac k2\log k$
and
$c_+=c+\tfrac k2\log\tfrac{2\pi e}k-h(U)$, where $c$ is the constant of
Theorem~\ref{thm:A}(d).
\end{corollary}

\begin{proof}
By Proposition~\ref{prop:B}(c), $\RC=R_U$. The Shannon lower bound gives
$R_U(D)\ge h(U)-\tfrac k2\log(2\pi eD/k)$. A Gaussian source has the largest
rate--distortion function among sources of a given covariance, and for $D\le k\kappa$
the Gaussian value is $\tfrac12\log\det K_U+\tfrac k2\log(k/D)$
\cite{coverthomas}. Combine these with Theorem~\ref{thm:A}(c) and (d).
\end{proof}

The corollary is stated with the output dimension $k$. When $J_g$ has rank $k$ at almost
every point, $k$ is the pointwise rank $r$, and with $q=p$ the gap is
$\tfrac{q-r}2\log(1/D)+O(1)$. The hypotheses on the law of $U$ are assumptions of the
corollary. Bounded derivatives of $g$ do not imply them.

\begin{remark}[A fixed kernel]
When $q=r$ the kernel does not move, and T1's remark on read operators of lower rank
treats that case under its sandwich assumption. Theorem~\ref{thm:A} needs $q=p$ and says
nothing about it when $r<p$. Theorem~\ref{thm:reduce} below does: it gives $\Rfix$ for every
$q\ge1$, and under the hypotheses of Corollary~\ref{cor:gapq} with $k=q$ the gap
$\Rfix-\RC$ is bounded, with no sandwich assumption. That is a statement about the
fixed-covariance class. It is not T1's bounded price for codes, which rests on the lower
quadratic sandwich, and a fixed kernel alone does not give that sandwich.
\end{remark}

\section{The constant, and the case of lower rank}\label{sec:sharp}

The estimates behind Theorem~\ref{thm:A} give the constant as well as the coefficient.
Theorems~\ref{thm:sharp} and~\ref{thm:reduce} are the main results of this note;
Theorem~\ref{thm:A} is the step on which they stand.
Write
\[
R_0(D)=\tfrac p2\log\frac pD+\tfrac12\log\det(\Sigma Q),\qquad B=K(\tr Q^{-1})^2 .
\]

\begin{theorem}[The fixed-covariance rate to within $O(D)$]\label{thm:sharp}
Suppose $Q\succ0$, with $\lambda$, $K$ and $D_0$ as in Theorem~\ref{thm:A}.
\begin{enumerate}
\item[(a)] For $D<D_0$,
$\Rfix(D)\ge R_0(D)-\tfrac p2\log\bigl(1+4KD/\lambda^2\bigr)$.
\item[(b)] For $D<\min\{p\,\lmin(\Sigma)\lambda,\ p^2/B\}$,
$\Rfix(D)\le R_0(D)-\tfrac p2\log\bigl(1-BD/p^2\bigr)$.
\end{enumerate}
Hence $\Rfix(D)=\tfrac p2\log\tfrac pD+\tfrac12\log\det(\Sigma Q)+O(D)$ as $D\to0$.
\end{theorem}

\begin{proof}
(a) Let $D<D_0$ and $L(\Delta)\le D$, and write $s=\tr\Delta$. By
Theorem~\ref{thm:A}(b), $s<s_0$, so by Theorem~\ref{thm:A}(a), $s\le2D/\lambda$. Then
$\tr(Q\Delta)\le L(\Delta)+Ks^2\le D(1+4KD/\lambda^2)$. By the inequality of the
arithmetic and geometric means applied to the eigenvalues of
$Q^{1/2}\Delta Q^{1/2}$, $\det(Q\Delta)\le(\tr(Q\Delta)/p)^p$. So
\[
r(\Delta)=\tfrac12\log\det(\Sigma Q)-\tfrac12\log\det(Q\Delta)
\ \ge\ \tfrac12\log\det(\Sigma Q)+\tfrac p2\log\frac{p}{D(1+4KD/\lambda^2)} .
\]
(b) Let $t>0$ solve $pt+Bt^2=D$ and take $\Delta=tQ^{-1}$. Then $t\le D/p$, so
$\Delta\preceq(t/\lambda)I\prec\lmin(\Sigma)I\preceq\Sigma$, and $\Delta$ is admissible.
Its trace is $t\tr Q^{-1}$, so
$L(\Delta)\le\tr(Q\Delta)+K(\tr\Delta)^2=pt+Bt^2=D$. Its rate is
$\tfrac12\log\det(\Sigma Q)+\tfrac p2\log(1/t)$, and
$pt=D-Bt^2\ge D(1-BD/p^2)>0$.
\end{proof}

The next theorem removes the hypothesis $Q\succ0$. Let $q=\operatorname{rank}Q\ge1$ and
$H=\Sigma^{1/2}Q\Sigma^{1/2}$, which also has rank $q$. Let $O$ be an orthogonal matrix
whose first $q$ columns, forming $O_1$, span the range of $H$. Define the reduced
consumer and source
\[
f(s)=g(\Sigma^{1/2}O_1s),\quad s\in\R^q,\qquad S\sim N(0,I_q),
\]
and write $\Rfix^{f}$ for the fixed-covariance rate of the pair $(S,f)$, defined as in
Section~\ref{sec:setting} with $q$, $I_q$ and $f$ in place of $p$, $\Sigma$ and $g$.

\begin{theorem}[Reduction to the read directions]\label{thm:reduce}
\begin{enumerate}
\item[(a)] $g(\Sigma^{1/2}Ow)=f(w_1,\dots,w_q)$ for every $w\in\R^p$.
\item[(b)] $\Rfix(D)=\Rfix^{f}(D)$ for every $D>0$.
\item[(c)] $f$ is $C^3$ with $\sup\|D^jf\|\le M_j\|\Sigma\|^{j/2}$, and
$Q_f=\E\,J_f(S)^{\mathsf T}J_f(S)=O_1^{\mathsf T}HO_1\succ0$, whose eigenvalues are the
nonzero eigenvalues of $H$.
\end{enumerate}
Hence Theorems~\ref{thm:A} and~\ref{thm:sharp} apply to $(S,f)$, and with $\det^{+}$ the
product of the nonzero eigenvalues,
\[
\Rfix(D)=\tfrac q2\log\frac qD+\tfrac12\log{\det}^{+}\bigl(\Sigma^{1/2}Q\Sigma^{1/2}\bigr)
+O(D)\qquad(D\to0).
\]
\end{theorem}

\begin{proof}
(a) Let $F(w)=g(\Sigma^{1/2}Ow)$ and $W=O^{\mathsf T}\Sigma^{-1/2}T\sim N(0,I_p)$, so that
$F(W)=g(T)$ and $J_F=J_g\Sigma^{1/2}O$. Then
$\E\,J_F(W)^{\mathsf T}J_F(W)=O^{\mathsf T}HO=\operatorname{diag}(Q_f,0)$, since the last
$p-q$ columns of $O$ span $\ker H$. The argument of Lemma~\ref{lem:kernel} applies to $F$
and $W$: for $j>q$, $\E\|J_F(W)e_j\|^2=0$, so $\partial_jF(W)=0$ almost surely. The
derivative is continuous and $W$ has full support, so $\partial_jF=0$ everywhere. Hence
$F$ does not depend on the last $p-q$ coordinates, and $F(w)=F(w_1,\dots,w_q,0)=f(w_1,\dots,w_q)$.

(b) The map $\Delta\mapsto\Delta'=O^{\mathsf T}\Sigma^{-1/2}\Delta\Sigma^{-1/2}O$ is a
bijection from $\{0\prec\Delta\preceq\Sigma\}$ onto $\{0\prec\Delta'\preceq I_p\}$, and
$r(\Delta)=-\tfrac12\log\det\Delta'$. In the channel of $\Delta$, put
$\hat W=O^{\mathsf T}\Sigma^{-1/2}\hat T$ and $Z'=O^{\mathsf T}\Sigma^{-1/2}Z$. These are
independent, with laws $N(0,I_p-\Delta')$ and $N(0,\Delta')$. Let $E$ be the leading
$q\times q$ block of $\Delta'$; then $0\prec E\preceq I_q$. By (a) the loss is
$\|f(\hat W_{1:q}+Z'_{1:q})-f(\hat W_{1:q})\|^2$, and $\hat W_{1:q}$ and $Z'_{1:q}$ are
independent with laws $N(0,I_q-E)$ and $N(0,E)$. So $L(\Delta)$ is the loss of the reduced
channel with error covariance $E$. For the rate, $\det\Delta'$ is $\det E$ times the
determinant of a Schur complement, which is positive definite and at most the trailing
block of $\Delta'$, hence at most $I_{p-q}$. So $\det\Delta'\le\det E$ and
$r(\Delta)\ge-\tfrac12\log\det E$. This gives $\Rfix\ge\Rfix^{f}$. Conversely, for any
$0\prec E\preceq I_q$ the choice $\Delta'=\operatorname{diag}(E,I_{p-q})$ has the same
loss and rate $-\tfrac12\log\det E$. It leaves the last coordinates undescribed, which is
allowed because the definition asks only $0\prec\Delta\preceq\Sigma$, so the reproduction
covariance $\Sigma-\Delta$ may be singular. If both problems are posed with the strict inequalities $\Delta\prec\Sigma$ and $E\prec I_q$
instead, the equality still holds: take
$\operatorname{diag}(E,(1-\varepsilon)I_{p-q})$, which has the same loss, since the loss
depends on $\Delta'$ only through $E$, and rate larger by
$\tfrac{p-q}2\log\tfrac1{1-\varepsilon}$, and let $\varepsilon\to0$.

(c) $J_f=J_g\Sigma^{1/2}O_1$ and the higher derivatives are bounded in the same way. The
computation in (a) gives $Q_f$, and $O_1^{\mathsf T}HO_1$ is $H$ restricted to its range.
For the last display, $\det(I_qQ_f)={\det}^{+}H$.
\end{proof}

No invariance of $\ker Q$ under $\Sigma$ is needed. Theorems~\ref{thm:A}
and~\ref{thm:sharp} are applied to $(S,f)$ with $q$, $I_q$ and $Q_f$ in place of $p$,
$\Sigma$ and $Q$; with $\lambda_f$, the least nonzero eigenvalue of $H$, in place of
$\lambda$; with $M_j\|\Sigma\|^{j/2}$ in place of $M_j$ and $q$ in place of $\tr\Sigma$ in
\eqref{eq:K}, which gives a value, and $K_f>0$ is any constant at least that value; and
with the $D_0$ of the pair $(S,f)$. The lower bound
then holds for $D$ below that $D_0$, and the upper bound for
$D<\min\{q\lambda_f,\ q^2/(K_f(\tr Q_f^{-1})^2)\}$. If $q=0$ the consumer is constant
and $\Rfix=0$.

\begin{corollary}[The rate gap in general rank]\label{cor:gapq}
Let $q=\operatorname{rank}Q\ge1$. Suppose the hypotheses of
Proposition~\ref{prop:B}(c) hold, that $U$ has a density with finite differential entropy
$h(U)$, and that its covariance $K_U$ is finite and positive definite. Then as $D\to0$,
\[
\tfrac{q-k}2\log\frac1D+c'_-+O(D)\ \le\ \Rfix(D)-\RC(D)\ \le\
\tfrac{q-k}2\log\frac1D+c'_++O(D),
\]
with $c'_-=\tfrac q2\log q+\tfrac12\log{\det}^{+}H-\tfrac12\log\det K_U-\tfrac k2\log k$
and $c'_+=\tfrac q2\log q+\tfrac12\log{\det}^{+}H+\tfrac k2\log\tfrac{2\pi e}k-h(U)$.
\end{corollary}

\begin{proof}
Theorem~\ref{thm:reduce} gives $\Rfix$, Proposition~\ref{prop:B}(c) gives $\RC=R_U$, and
$R_U$ lies between the Shannon lower bound and the Gaussian upper bound, as in the proof
of Corollary~\ref{cor:C}.
\end{proof}

When $J_g$ has rank $k$ almost everywhere, so that $r=k$, the coefficient is $(q-r)/2$.
The corollary does not cover an output whose law is singular, as on a curved image of
dimension less than $k$. The two constants differ by
$\tfrac12\log((2\pi e)^k\det K_U)-h(U)$, which is zero when $U$ is Gaussian.

\section{An example}\label{sec:example}

Let $p=2$, $\Sigma=I$, and
\[
g_a(x_1,x_2)=x_1+a\sin x_2,\qquad a\ne0 .
\]
Its derivatives are bounded: $M_1\le(1+a^2)^{1/2}$ and $M_2,M_3\le|a|$. The Jacobian
$(1,\ a\cos x_2)$ has rank one everywhere and its kernel turns with $x_2$. With
$\E\cos X_2=e^{-1/2}$ and $\E\cos^2X_2=\tfrac12(1+e^{-2})$,
\[
Q=\begin{pmatrix}1&ae^{-1/2}\\ ae^{-1/2}&\tfrac{a^2}2(1+e^{-2})\end{pmatrix},\qquad
\det Q=\tfrac{a^2}2(1-e^{-1})^2>0 .
\]
So $q=2$ and $r=k=1$. The image is $\R$, which is closed and convex, and $s(u)=(u,0)$
is a section. The output $U=X_1+a\sin X_2$ is a Gaussian variable plus an independent
bounded one, so it has a smooth density, finite differential entropy, and variance
$1+\tfrac{a^2}2(1-e^{-2})$. Corollary~\ref{cor:C} applies and the gap is
$\tfrac12\log(1/D)+O(1)$; Proposition~\ref{prop:gapex} below sharpens this to
$\tfrac12\log(1/D)+C_*+O(D)$ with an explicit constant $C_*$. At $a=1$, $h(U)=1.598216$ nats. This is a quadrature, not a
Monte Carlo estimate: the density of $U$ is the one-dimensional integral
$f(u)=\E\,\varphi(u-\sin X_2)$ with $\varphi$ the standard normal density, and
$h(U)=-\int f\log f$. Two resolutions in the exact script described below (Gauss--Hermite rules of
80 and 240 points for $f$; Simpson's rule on 4{,}001 and 32{,}001 points for the outer
integral) agree to twelve decimals, and the same quadrature returns the variance
$1+\tfrac12(1-e^{-2})$ to twelve decimals. It is a numerical value with no proved error
bound.

\begin{proposition}[The loss in closed form, and a range for the lower bound]\label{prop:exact}
Let $\Sigma=I_2$ and $0\preceq\Delta\preceq I_2$ with entries $d_{11}$, $d_{12}$ and
$d_{22}=v$. Let
\[
F(v)=1-e^{-v/2}+e^{-2}\Bigl(e^{3v/2}-\tfrac12(1+e^{2v})\Bigr).
\]
\begin{enumerate}
\item[(a)] $L(\Delta)=d_{11}+2ae^{-1/2}d_{12}+a^2F(v)$.
\item[(b)] $F$ is concave on $[0,1]$ with $F(0)=0$ and $F(1)=\tfrac12(1-e^{-2})$, so
$L(\Delta)\ge\tr(\Theta_a\Delta)\ge\mu\tr\Delta$, where
\[
\Theta_a=\begin{pmatrix}1&ae^{-1/2}\\ ae^{-1/2}&\tfrac{a^2}2(1-e^{-2})\end{pmatrix},\qquad
\det \Theta_a=a^2\Bigl(\tfrac12(1-e^{-2})-e^{-1}\Bigr)>0,
\]
and $\mu=\lmin(\Theta_a)>0$.
\item[(c)] $D_0\ge\mu s_0$.
\end{enumerate}
\end{proposition}

\begin{proof}
(a) Write $A$ and $B$ for the second coordinates of $\hat T$ and $Z$; they are independent
with variances $1-v$ and $v$. The loss is $(Z_1+a(\sin(A+B)-\sin A))^2$. Its first term
has mean $d_{11}$. For the cross term, $Z_1$ is independent of $A$, and by Gaussian
integration by parts $\E[Z_1\sin(A+B)]=d_{12}\,\E\cos(A+B)=d_{12}e^{-1/2}$. For the last
term, with $\E\cos N=e^{-\sigma^2/2}$ for $N\sim N(0,\sigma^2)$:
$\E\sin^2(A+B)=\tfrac12(1-e^{-2})$, $\E\sin^2A=\tfrac12(1-e^{-2+2v})$, and
$\E[\sin(A+B)\sin A]=\tfrac12(\E\cos B-\E\cos(2A+B))=\tfrac12(e^{-v/2}-e^{-2+3v/2})$.
Adding gives $a^2F(v)$.

(b) With $y=e^{v/2}\ge1$, $4e^{v/2}F''(v)=-1+e^{-2}y^4(9-8y)$, and $y^4(9-8y)$ decreases
for $y\ge1$ from the value $1$, so $F''<0$. A concave function with $F(0)=0$ lies above its
chord on $[0,1]$: $F(v)\ge vF(1)$. So $L(\Delta)\ge d_{11}+2ae^{-1/2}d_{12}+a^2F(1)v
=\tr(\Theta_a\Delta)\ge\mu\tr\Delta$, as $\Delta\succeq0$.

(c) If $\tr\Delta\ge s_0$ then $L(\Delta)\ge\mu s_0$.
\end{proof}

At $a=1$: $\mu=0.04651$, and with $s_0=0.0100$ the range of Theorem~\ref{thm:A}(c)
contains every $D<\mu s_0=0.000465$. With Remark~\ref{rem:scope},
$0.000465\le D_0\le0.0021$. In the notation of Corollary~\ref{cor:gapq},
$c'_-=-0.291754$ and $c'_+=-0.291379$. These differ by a fixed $0.0004$, so the corollary
alone places the gap between two expansions and does not give one constant. For this
example a further step does.

\begin{proposition}[The gap for the example]\label{prop:gapex}
For $g_a$ with $\Sigma=I_2$, let
$\mathrm{SLB}(D)=h(U)-\tfrac12\log(2\pi eD)$. Then for every $D>0$,
\[
\mathrm{SLB}(D)\ \le\ \RC(D)\ \le\ \mathrm{SLB}(D)+\tfrac12\log(1+D),
\]
and therefore, as $D\to0$,
\[
\Rfix(D)-\RC(D)=\tfrac12\log\frac1D+C_*+O(D),\qquad
C_*=\log2+\tfrac12\log\det Q+\tfrac12\log(2\pi e)-h(U).
\]
\end{proposition}

\begin{proof}
By Proposition~\ref{prop:B}(c), $\RC=R_U$, and the lower bound is the Shannon lower bound
for the scalar $U$. For the upper bound write $U=Y+G$ with $Y=a\sin X_2$ and $G=X_1$,
which are independent, and take $\hat U=U+\sqrt D\,N$ with $N$ standard normal and
independent of $U$. Then $\E(\hat U-U)^2=D$ and
$I(U;\hat U)=h(U+\sqrt D\,N)-\tfrac12\log(2\pi eD)$. For $\sigma^2\ge1$ let $G_\sigma$ be
normal with variance $\sigma^2$ and independent of $Y$; then
$h(Y+G_\sigma)=I(Y;Y+G_\sigma)+\tfrac12\log(2\pi e\sigma^2)$. The mutual information does
not increase with $\sigma$, by data processing, since $Y+G_{\sigma'}$ is $Y+G_\sigma$ plus
independent noise when $\sigma'>\sigma$. So
$h(U+\sqrt D\,N)-h(U)\le\tfrac12\log(1+D)$, which gives the upper bound. For the
expansion, $R_0(D)-\mathrm{SLB}(D)=\tfrac12\log(1/D)+C_*$ exactly, and
Theorem~\ref{thm:sharp} with $p=2$ gives both sides. For $D<D_0$, part (a) and the upper
bound on $\RC$ give
$\Rfix-\RC\ge\tfrac12\log\tfrac1D+C_*-\log(1+4KD/\lambda^2)-\tfrac12\log(1+D)$, and
$D_0\ge\mu s_0>0$ by Proposition~\ref{prop:exact}(c). For $D<\min\{2\lambda,\ 4/B\}$, part
(b) and the lower bound on $\RC$ give
$\Rfix-\RC\le\tfrac12\log\tfrac1D+C_*-\log(1-BD/4)$.
\end{proof}

$C_*$ is the constant $c'_+$ of Corollary~\ref{cor:gapq}: for this consumer the Shannon
lower bound is tight to $O(D)$, and the Gaussian upper bound on $R_U$ is not. At $a=1$,
$C_*=-0.291379$, with $h(U)$ from the quadrature above.

\begin{sloppypar}
\textbf{Numerical check with the exact loss.} The script
\texttt{verify\_ft7\_exact.py} in \texttt{experiments/} (seed 20261122, output
\texttt{FT7-exact-net.json} in \texttt{results/}) uses the closed form of
Proposition~\ref{prop:exact}. The closed form agrees with Monte Carlo within $1.2$
standard errors in six cases: three covariances at each of $a=1$ and $a=0.7$. On 20{,}000 random admissible
$\Delta$, $L\ge\tr(\Theta_a\Delta)$ always, and the least sampled ratio $L/\tr\Delta$ is $0.083$,
against $\mu=0.0465$; the constant $\mu$ is a valid bound and is not sharp. A numerical
estimate of $D_0$ is $0.0014$, inside $[0.000465,0.0021]$; it is a search, not a certified
infimum. With the exact loss, a search for the least rate over all $2\times2$ covariances
gives $9.09817$, $11.40082$ and $13.70341$ nats at $D=10^{-4}$, $10^{-5}$ and $10^{-6}$,
inside the bounds of Theorem~\ref{thm:sharp}; it is below $R_0(D)$ by $0.68D$ at each. The
script also evaluates the family $\Delta=tQ^{-1}$ with the exact loss: at $D=10^{-5}$ its
rate is $11.400817$ nats, and the search's least rate agrees with it to $10^{-9}$. At $D=10^{-5}$ the family's rate exceeds the Shannon lower bound on $R_U$ by $5.465077$
nats, against $\tfrac12\log(1/D)+C_*=5.465084$ from Proposition~\ref{prop:gapex}; the
difference is $7\cdot10^{-6}$, the size of the $O(D)$ term. The value of $h(U)$ cancels in
this comparison, so it tests the constant of Theorem~\ref{thm:sharp} and not $h(U)$; the
script checks $h(U)$ separately, by a second quadrature method. It exceeds the Gaussian upper
bound on $R_U$ by $5.464702$, which is $0.0004$ less because that bound is not tight. The
script also checks the entropy step of Proposition~\ref{prop:gapex} by quadrature at
$D=0.1$ and $0.01$. The search is not a certified minimum. For
Theorem~\ref{thm:reduce} the script takes $p=3$, a correlated $\Sigma$, and
$g(x)=g_1(Bx)$ with $B\Sigma B^{\mathsf T}=I_2$, for which $q=2$ and $\ker Q$ is not
invariant under $\Sigma$. The loss of the three-dimensional channel, by Monte Carlo,
agrees with the closed form at the block $E=B\Delta B^{\mathsf T}$ within $1.5$ standard
errors on four random $\Delta$; the rate exceeds $-\tfrac12\log\det E$ on each; and the
lift attains it. This consumer is $g_1$ composed with a linear map, so part (a) of the
theorem holds by construction and only part (b) is tested. The fire test has six mutants,
each a changed formula.
\end{sloppypar}

\textbf{An earlier check by Monte Carlo.} The script
\texttt{experiments/verify\_ft7\_moving\_kernel.py} (seed 20261131, output
\texttt{results/FT7-moving-kernel-net.json}) predates the closed form. It works at $a=1$ by Monte Carlo with common random numbers, so its figures carry
sampling error that the exact check above does not. It is kept because it also covers
other consumers. It makes three kinds of check, and they test different things.

\emph{The lower bound.} For $g_1$ the script searches numerically for the least rate
over all $2\times2$ covariances $\Delta$ (rotation, anisotropy and scale; a grid and a
simplex polish, not a certified minimum) subject to the loss constraint. At
$D=10^{-2},10^{-3},10^{-4},10^{-5}$ it finds $4.485$, $6.794$, $9.097$ and
$11.400$ nats, with slope $1.0009$ in $\log(1/D)$ against $p/2=1$. These sit above the
expression of Theorem~\ref{thm:A}(c), which with the exact $\lambda=0.13993$ is $2.639$,
$4.941$, $7.244$ and $9.546$. Part (c) is proved only for $D<D_0$. By
Remark~\ref{rem:scope}, $D_0\le3\lambda^2/(4K)=0.0021$ here, so $D=10^{-2}$ is outside
the proved range. By Proposition~\ref{prop:exact}, $D_0\ge0.000465$, so $D=10^{-4}$ and $10^{-5}$ are
inside it. $D=10^{-3}$ is not shown to be inside. The
gap to that expression is a constant, about $1.85$ nats, so the slope is the sensitive
check. The inequality of Theorem~\ref{thm:A}(a) is checked on random $\Delta$ and on
$\Delta$ nearly of rank one along the least eigenvector of $Q$, where it is closest to
equality: the ratio $\lambda\tr\Delta/(2L)$ tends to $1/2$ there, and is between $0.494$
and $0.498$ on those cases, against the bound $1$. Here $K=6.993$ and $s_0=0.0100$, and
most of the sampled range $\tr\Delta\in[10^{-4},0.03]$ lies beyond $s_0$, so the check
reaches past what part (a) proves.

\emph{The upper bound.} The family $\Delta=tQ^{-1}$ of part (d), evaluated on the same
sample as the search, agrees with the search's minimum to within $10^{-4}$ nats at all
four values of $D$. On a separate, larger sample it gives $4.489$, $6.798$, $9.101$ and
$11.403$ nats; the difference of $0.004$ is Monte Carlo noise between the two samples. Over
$D\in[10^{-5},10^{-2}]$ the family's slope is $1.0007$. At $D=10^{-5}$ the family's rate on the larger
sample exceeds $R_U$ by between $5.4672$ and $5.4676$ nats, the two ends coming from the
Gaussian upper bound and the Shannon lower bound on $R_U$. Those two bounds on $R_U$ are
$0.0004$ apart, while the family's rate differs by $0.004$ between the two samples, so
the sampling error is the larger uncertainty and the gap figures are good to about two
decimals: the family's rate exceeds $R_U$ by about $5.46$ to $5.47$ nats. The family's
rate is an upper bound on $\Rfix$, so this bounds the gap $\Rfix-\RC$ from above, up to
that sampling error, against $\tfrac12\log(1/D)+c_+=6.158$. It is the comparison of part
(d) again, not a separate test of Corollary~\ref{cor:C}. From below, the search's minimum
(on the smaller sample) less the Gaussian upper bound on $R_U$ is $5.46$ nats, against
$\tfrac12\log(1/D)+c_-=3.610$, a bound that holds at this $D$, since $10^{-5}<0.000465\le D_0$ by
Proposition~\ref{prop:exact}.
That is the comparison of part (c) again, and it rests on a search that is not a
certified minimum. What the corollary adds to Theorem~\ref{thm:A}, namely $\RC=R_U$ and
the two bounds on $R_U$, has no numerical check beyond the identity $g(s(u))=u$, which is
trivial for this $g$.

\emph{Other consumers, upper bounds only.} For three smooth radial consumers the script
evaluates one hand-picked family each, and finds slopes $1.0002$, $1.4995$ and $1.5001$
against $1$, $1.5$ and $1.5$. These are rates of achievable families. They test part
(d) and T1's second-order statement, not the lower bound. A control with a fixed kernel,
$g(x)=x_1+0.3\sin x_1$, has slope $0.5000$ and an excess between $0.004$ and $0.014$
nats; for a control an upper bound is the right direction.

The script's fire test has six mutants. Two run the full search on a changed instrument:
a fixed-kernel consumer, and a loss that reads one fixed direction at $37^\circ$, off the
search's angle grid. Both must return slope $1/2$ in place of $1$, and both do. The
other four change a predicted number.

\section{What is not claimed}

\textbf{Codes.} Every statement about $\Rfix$ is about fixed-covariance Gaussian test
channels. It is not a lower bound for every code built from a weight. A test channel
that is not of fixed covariance, for instance one that describes $U=g(T)$ and reproduces
through a right inverse, can have a lower rate; Proposition~\ref{prop:B} and the
corollaries compare information rates of that kind. No statement about codes is made.

\textbf{A model case.} Three bounded derivatives of $g$ and a Gaussian source are used
throughout. They exclude natural consumers such as $g(x)=\|x\|^2$. The results are a model
case, not a general theory of non-difference distortion measures.

\textbf{The range of $D$.} Theorem~\ref{thm:A}(c) and Theorem~\ref{thm:sharp}(a) hold for
$D<D_0$, and in general we have no explicit lower bound for $D_0$. The sine example has
one, from a concavity special to it. So for a general consumer the lower bounds are
asymptotic statements with no explicit range of $D$.

\textbf{Learning.} Proposition~\ref{prop:B} assumes the decoder holds a section. How a
learner obtains one from queries, and at what cost, depends on the query model: the
noise, the precision, the calibration message, and the number of queries per served
symbol. For the family $g_a$ with only $a$ unknown, one exact query at $(0,\pi/2)$
returns $a$, whatever $D$ is. So any growth of the query cost as $D\to0$ has to come
from noise, limited precision, or the size of the consumer class. That question is a separate item of the program and is not treated here. Note also that
the right inverse $s(u)=(u,0)$ does not depend on $a$: obtaining a right inverse and
identifying the whole consumer need not be the same task. Obtaining a right inverse means
learning a selection of fiber representatives, one for each level set, which is not
the same as learning $Q$. Until that problem is treated, the price in this note is the
price of stopping at $Q$ within the fixed-covariance class. It is not a price of learning
$g$.

\textbf{Output laws and images.} Corollaries~\ref{cor:C} and~\ref{cor:gapq} need a
closed convex image and an output with a density in $\R^k$, finite differential entropy
and a positive definite covariance. That leaves out much of the geometry of interest. The
map $g(x)=(\cos x,\sin x)$ has a circle for its image, which is not convex, and
$g(x)=(x_1,\sin x_1)$ has a sine curve, on which the output has no density in $\R^2$.
Both maps have bounded derivatives and a measurable right inverse. For such
consumers Proposition~\ref{prop:B}(b) still gives $\RC=R_U^{\cU}$ when a right inverse
exists, but we do not evaluate $R_U^{\cU}$. A natural next step is to stratify the source
space by the rank of $J_g$, treat the locus where the rank drops, and ask what
$R_U^{\cU}$ is on each stratum.

\textbf{Prior work.} The best code for a locally quadratic loss, as opposed to a constant
weight, is the subject of the high-resolution theory of non-difference distortion
measures \cite{linderzamir,linderzamirzeger}. Quantization judged by the error of a
downstream function is studied as functional quantization by Misra, Goyal and Varshney
\cite{misra}. A review of v0.2 of this note pointed to the passages below. The drafter
then read them in the author-hosted preprints of the first two papers (not the published
versions) and the author copy of the published third, on 22 November 2026: the passages named, not the papers in full.

The main theorem of \cite{linderzamir} (Theorem~1, Section~2.1) gives the low-distortion
form of the rate--distortion function for a distortion measure with sensitivity matrix
$M(x)$. Its condition (b) asks that $d(x,y)$ have a unique minimum in $y$; condition (e),
$\E|\log\det M(X)|<\infty$, makes $M(X)$ positive definite almost surely, as the paper
notes; condition (g) bounds a moment of $\tr M^{-1}(X)$. Its concluding remarks name, as
problems for further research, a minimizer that is not piecewise invertible, source and
reproduction of different dimensions, and the case in which $M(X)$ is singular with
positive probability. The companding paper \cite{linderzamirzeger} (Section~3.1,
conditions (b) and (c)) assumes that $d(x,y)=0$ only at $y=x$ and that $M(x)$ is positive
definite for all $x$. The consumers of interest here, those with $k<p$, have a read metric $\PC(x)$ that is
singular at every point, and many reproductions of zero loss, so they fall outside those
hypotheses. That is a difference from those two theorems. It is not evidence of priority
over later work, which has not been surveyed.

In \cite{misra}, Section~I-B remarks that if the function were known to the encoder,
``one could do no better than to simply compute the function and encode the result.''
That is the idea of Proposition~\ref{prop:B}. What the proposition adds, if anything, is
the reproduction constraint made explicit through a right inverse and the image. Section~VII
of \cite{misra} treats intervals on which the function's sensitivity is zero. Its
Example~6, the function $\min(X,\tfrac12)$ of a uniform variable, has distortion
proportional to $2^{-2R}$ at fixed rate and to $2^{-4R}$ at variable rate, and the paper
notes that the exponent is larger. So the gains of functional quantization are not all
constant offsets, and a change of exponent caused by what the function ignores is not new
in itself. The setting here differs in two ways. In the example here the pointwise read rank is one everywhere and the
ignored direction turns with the point, where their example, which also has a scalar
output throughout, has a region where the local rank drops to zero. And the comparison here is against one explicitly restricted class, the
fixed-covariance Gaussian test channels, where theirs is between fixed-rate and
variable-rate quantizers. We did not read the proof of their Theorem~16 or the rest of the
paper. This note makes no claim of priority. What it adds to T1 is the comparison: under
the hypotheses of Corollary~\ref{cor:gapq} with $k<q$, the fixed-covariance Gaussian test
channels behind the codes T1 serves exceed $\RC$ by $\tfrac{q-k}2\log(1/D)+O(1)$.

\begin{thebibliography}{9}
\bibitem{t1} A. H. Bond, ``Learn then serve: Recovering a consumer's read operator and
paying for it'' (T1, working note, v4.1, 2026-11-18),
\texttt{paper/t1-learn-then-serve}.
\bibitem{coverthomas} T. M. Cover and J. A. Thomas, \emph{Elements of Information
Theory}, 2nd ed. Wiley, 2006, ch.~10.
% verify the chapter's statements of the Shannon lower bound and the Gaussian upper bound on final
\bibitem{linderzamir} T. Linder and R. Zamir, ``High-resolution source coding
for non-difference distortion measures: The rate-distortion function,''
\emph{IEEE Trans. Inf. Theory}, vol.~45, no.~2, pp.~533--547, 1999.
% Bibliographic data from DBLP. Section 2.1 and the concluding remarks read in the author-hosted version, 2026-11-22.
\bibitem{linderzamirzeger} T. Linder, R. Zamir, and K. Zeger, ``High-resolution source
coding for non-difference distortion measures: Multidimensional companding,''
\emph{IEEE Trans. Inf. Theory}, vol.~45, no.~2, pp.~548--561, 1999.
% Pages from a search result's reading of the PDF header. Section 3.1 read in the author-hosted version, 2026-11-22.
\bibitem{misra} V. Misra, V. K. Goyal, and L. R. Varshney, ``Distributed scalar
quantization for computing: High-resolution analysis and extensions,''
\emph{IEEE Trans. Inf. Theory}, vol.~57, no.~8, pp.~5298--5325, 2011.
% Bibliographic data from a web search (BU and MIT repository records, arXiv:0811.3617). Section I-B and Section VII through Theorem 15 read in the author copy of the published version, 2026-11-22.
\end{thebibliography}

\end{document}
